\documentclass[12pt,a4paper]{article}
\usepackage[french]{babel}
\usepackage[T1]{fontenc}
\usepackage[utf8]{inputenc}
\usepackage{color}
\usepackage{fancyhdr}
\usepackage{graphicx}
\usepackage{tcolorbox}
\usepackage{tabularx}
\usepackage{needspace}
\setlength{\headheight}{15pt}
\pagestyle{fancy}
\begin{document}

    \begin{figure}[t]
        \includegraphics[scale=0.5]{assets/logo\_lmu.png}
        \hfill
        \includegraphics[scale=0.5]{assets/logo\_ic2.png}
    \end{figure}

    \title{
        \textcolor{blue}{Le Mans Université}\\
        Licence Informatique \textit{2ème année}\\
        Module 174UP02 Rapport de Projet\\
        \textbf{Codenames}
    }
    \author{PIAU Romain - ROGER Noam \\QUINTON Chloé - MAUDET Mathis}
    \date{\today}
    \maketitle

    \newpage

    \tableofcontents

    \newpage 

    \section{Introduction}
        Ce document présente le projet Codenames, développé dans le cadre du projet de 2ème année de licence informatique à l'Université du Mans.
        Le projet consiste en la création d'un jeu numérique basé sur le concept du jeu de société Codenames.
        Les fonctionnalités principales incluent la gestion des parties, l'interface utilisateur et l'intégration des règles du jeu. 
        Ce rapport détaille les choix techniques, l'architecture du système et les résultats obtenus lors des tests.

    \section{Conception}

        \subsection{Présentation du jeu}

            \textit{Codenames} est un jeu d’association d’idées en équipes.
            Chaque équipe doit retrouver ses mots secrets à partir d’indices donnés par son maître-espion.

            \begin{tcolorbox}[colback=blue!5,colframe=blue!50]
            \textbf{Objectif :} être la première équipe à retrouver tous ses agents.
            \end{tcolorbox}

            \begin{enumerate}
                \item Mise en place :

                    \begin{itemize}
                        \item Former 2 équipes : rouge et bleue
                        \item Choisir un maître-espion par équipe
                        \item Disposer 25 mots en grille (5x5)
                        \item Tirer une carte clé au hasard\\ \textbf{Important :} La carte clé indique :
                        \begin{itemize}
                            \item 9 mots pour une équipe (celle qui commence)
                            \item 8 mots pour l’autre
                            \item 7 civils
                            \item 1 assassin
                        \end{itemize}
                    \end{itemize}

                \item Déroulement d’un tour

                    \begin{enumerate}
                        \item Donner un indice

                            Le maître-espion doit dire :

                            \begin{itemize}
                                \item Un mot unique
                                \item Un nombre
                            \end{itemize}

                            \vspace{0.5cm}

                            \needspace{5\baselineskip}

                            \textbf{Exemple :}

                            \begin{tcolorbox}
                            Indice : \textbf{Animal 2}
                            \end{tcolorbox}

                            Cela signifie que 2 mots sur la table sont liés à "animal".

                            \vspace{0.5cm}
                            
                            \textbf{Règles strictes pour les indices}

                            \begin{itemize}
                                \item Un seul mot (pas de phrase)
                                \item Pas un mot présent sur la table
                                \item Pas une partie d’un mot (ex : "auto" pour "automobile" interdit)
                                \item Pas de traduction directe (ex : dire "dog" pour "chien")
                                \item Pas de communication non verbale
                            \end{itemize}


                        \item Phase de devinettes

                            Les joueurs peuvent proposer des mots.

                            \begin{itemize}
                                \item Nombre maximum de propositions = nombre annoncé + 1
                            \end{itemize}

                            \vspace{0.5cm}

                            \textbf{Exemple :}

                            \begin{tcolorbox}
                            Indice : \textbf{Animal 2}
                            \end{tcolorbox}

                            → L’équipe peut proposer jusqu’à 3 mots.

                            \vspace{0.5cm}

                            Résultat des choix

                            \begin{itemize}

                                \item Bon mot

                                    \begin{itemize}
                                        \item On recouvre avec la couleur de l’équipe
                                        \item L’équipe peut continuer
                                    \end{itemize}

                                \item Mot neutre

                                    \begin{itemize}
                                        \item Tour terminé immédiatement
                                    \end{itemize}

                                \item Mot adverse

                                    \begin{itemize}
                                        \item Tour terminé
                                        \item Aide l’équipe adverse
                                    \end{itemize}

                                \item Assassin

                                    \begin{tcolorbox}[colback=red!5,colframe=red!50]
                                    \textbf{Effet : défaite immédiate}
                                    \end{tcolorbox}
                                
                            \end{itemize}

                    \end{enumerate}
                
                \item Fin de partie

                    \begin{itemize}
                        \item Une équipe gagne si elle trouve tous ses mots
                        \item Une équipe perd immédiatement si elle choisit l’assassin
                    \end{itemize}

                \item Règles avancées

                    Indice avec 0

                    \begin{tcolorbox}
                    Exemple : \textbf{Animal 0}
                    \end{tcolorbox}

                    Signifie :
                    Aucun mot restant n’est lié à "animal".

                    → Permet d’éviter certains pièges.

            \end{enumerate}
        
        \Needspace{10\baselineskip}
        \subsection{Notre adaptation numérique}
            Notre version numérique de Codenames vise à recréer l'expérience du jeu de société tout en offrant des fonctionnalités supplémentaires pour améliorer l'interaction et la convivialité.

            Fonctionnalités principales :
            \begin{enumerate}
                \item Interface utilisateur intuitive grâce à SDL
                    \begin{itemize}
                        \item différentes intefaces créées pour chaque étape (menu, lobby, jeu),
                        \item destion avancée de bouttons et inputs,
                        \item bandeau d'informations (affichage fps, réglage du volume et de la luminosité, règles du jeu)...
                    \end{itemize}
                \item Gestion automatique des règles du jeu, le code permet  
                    \begin{itemize}
                        \item de vérifier la validité des indices donnés par les maîtres-espions,
                        \item de gérer les différentes phases du jeu (donner un indice, devinettes, fin de partie),
                        \item d'appliquer les conséquences des choix des joueurs (bon mot, mot neutre, mot adverse, assassin).
                    \end{itemize}
                \item Possibilité de jouer en ligne avec des amis grâce à la création de lobbies privés
                \item Système de chat intégré pour communiquer entre les joueurs
                \item Historique des indices 
            \end{enumerate}

    \newpage
    \section{Organisation du travail}
        \renewcommand{\arraystretch}{1.3}
        \setlength{\tabcolsep}{8pt}
        \begin{table}[h]
        \centering
        \begin{tabularx}{\textwidth}{|>{\raggedright\arraybackslash}X|c|c|c|c|}
        \hline
        \textbf{Tâches} & \textbf{Noam} & \textbf{Romain} & \textbf{Chloé} & \textbf{Mathis} \\
        \hline
        Conception de la structure du code & 40\% & 30\% & 15\% & 15\% \\
        \hline
        Fonctions principales (serveur) & 30\% & 30\% & 20\% & 20\% \\
        \hline
        Fonctions principales (client) & 20\% & 20\% & 30\% & 30\% \\
        \hline
        Affichage de l'interface graphique & 10\% & 30\% & 30\% & 30\% \\
        \hline
        Gestion du lobby & 20\% & 20\% & 15\% & 45\% \\
        \hline
        Gestion des fonctions du jeu & 25\% & 25\% & 25\% & 25\% \\
        \hline
        Rédaction du rapport de projet & 15\% & 15\% & 55\% & 15\% \\
        \hline
        Tests et débogage & 25\% & 25\% & 25\% & 25\% \\
        \hline
        \end{tabularx}
        \caption{Répartition des tâches entre les membres de l'équipe}
        \end{table}

    \newpage
    \section{Développement}

        \subsection{Structure du code}

        Le code du jeu est structuré de la manière suivante :

        \begin{itemize}
            \item Deux parties principales : client et serveur
            \item Un dossier de tests
            \item Un dossier de documentation
            \item Un dossier rapport 
            \item Un fichier README.md pour expliquer les principales fonctionnalités du jeu et les instructions d'installation et d'utilisation
            \item Un fichier TODO.md pour suivre les tâches à accomplir et les fonctionnalités à implémenter
            \item Un fichier .gitignore pour exclure les fichiers temporaires et les ressources non nécessaires du dépôt Git
            \item Un fichier LICENSE pour spécifier les conditions de distribution et d'utilisation du code
            \item Un fichier setup pour faciliter l'installation des dépendances et la configuration de l'environnement de développement
            \item Un fichier doxyfile pour la génération de la documentation avec Doxygen
        \end{itemize}

        Les parties client et serveur sont les plus développées et contiennent les fonctionnalités principales du jeu. Elles partagent une structure similaire pour faciliter la maintenance et la compréhension du code.
        \begin{enumerate}
            \item \textbf{Points communs des deux parties}
                \begin{itemize}
                    \item Dossier assets : contient les ressources graphiques, textuelles ou sonores du jeu
                    \item Dossier lib : contient les fichiers d’en-tête et les bibliothèques
                    \item Dossier src : contient le code source
                    \item Fichier Makefile : pour la compilation du projet
                    \item Fichier run : pour lancer le jeu
                \end{itemize}
            \item \textbf{Côté client}
                \begin{itemize}
                    \item Bibliothèques SDL compilés (ajouté par le script setup)
                    \item Gestion de l'interface graphique (menu, lobby, jeu) 
                    \item Traitement des bouttons, inputs utilisateur, crossfade, bandeau d'informations...
                    \item Communication avec le serveur (envoi des actions du joueur, réception des mises à jour du jeu)
                \end{itemize}
            \item \textbf{Côté serveur}
                \begin{itemize}
                    \item Gestion des connexions clients
                    \item Implémentation de la logique du jeu
                \end{itemize}
        \end{enumerate}

        \subsection{Principales structures de données}

        \begin{itemize}
            \item \textbf{AppContext} : contient les informations sur l’état de l’application.\\
                Utilité : centraliser les données et les ressources du jeu pour faciliter la gestion de l’état et de l’affichage.
                Détails :
                \begin{itemize}
                    \item \textbf{Gérer l’affichage SDL}
                    \begin{itemize}
                        \item \texttt{window} $\rightarrow$ la fenêtre du jeu
                        \item \texttt{renderer} $\rightarrow$ le moteur de rendu graphique
                        \item \texttt{bg\_color} $\rightarrow$ couleur de fond
                    \end{itemize}

                    \item \textbf{Gérer l’état du jeu}
                    \begin{itemize}
                        \item \texttt{app\_state} $\rightarrow$ menu, lobby, en partie, etc.
                        \item \texttt{lobby} $\rightarrow$ données du lobby (joueurs, salle, etc.)
                    \end{itemize}

                    \item \textbf{Gérer le joueur}
                    \begin{itemize}
                        \item \texttt{player\_role} $\rightarrow$ agent ou spy
                        \item \texttt{player\_team} $\rightarrow$ équipe
                        \item \texttt{player\_id} $\rightarrow$ identifiant serveur
                        \item \texttt{player\_uuid}, \texttt{player\_name} $\rightarrow$ identité persistée
                    \end{itemize}

                    \item \textbf{Gérer la connexion réseau}
                    \begin{itemize}
                        \item \texttt{sock} $\rightarrow$ socket réseau
                        \item \texttt{ping\_ms} $\rightarrow$ latence avec le serveur
                    \end{itemize}

                    \item \textbf{Gérer le timing et les performances}
                    \begin{itemize}
                        \item \texttt{clock} $\rightarrow$ nombre de frames écoulées
                        \item \texttt{fps} $\rightarrow$ images par seconde
                        \item \texttt{frame\_start\_time} $\rightarrow$ début de frame
                        \item \texttt{last\_fullscreen\_toggle\_ms} $\rightarrow$ évite les toggles rapides
                    \end{itemize}

                    \item \textbf{Gérer l’audio}
                    \begin{itemize}
                        \item \texttt{music\_volume}
                        \item \texttt{sound\_effects\_volume}
                    \end{itemize}

                    \item \textbf{Gérer la fenêtre}
                    \begin{itemize}
                        \item taille et position en mode fenêtré \\
                            (\texttt{windowed\_width}, \texttt{windowed\_height}, etc.)
                    \end{itemize}

                    \item \textbf{Informations générales}
                    \begin{itemize}
                        \item \texttt{version} $\rightarrow$ version du jeu
                    \end{itemize}
                \end{itemize}
        \end{itemize}

        \subsection{Interface graphique}

        \subsection{Gestion des parties}

        \subsection{Fonctionnalités additionnelles}

    \newpage
    \section{Bilan et conclusion}

    \newpage
    \section{Annexes}


\end{document}